
\section{Annihilation of positronium}

\SourcePage{413}{418}

By virtue of conservation of momentum the annihilation of the electron and positron in positronium must be accompanied by the emission of at least two photons. Such a decay is possible, however (principally) only for parapositronium; and for orthopositronium it is three-photon annihilation. In \S\,9 it was shown that the total angular momentum of a system of two photons cannot equal~1. Therefore an orthopositronium atom, which is in a state ${}^3S_1$, cannot decay into two photons. Moreover, since in the ${}^3S_1$ state positronium constitutes a charge-odd system (see the problem to \S\,27), by the Furry theorem (see \S\,79) its decay into any even number of photons is also impossible. On the contrary, in the ${}^1S_0$ state positronium is charge-even, and therefore the decay of parapositronium into any odd number of photons is forbidden.

The principal process determining the lifetime of positronium is, thus, two-photon annihilation in the case of parapositronium and three-photon annihilation in the case of orthopositronium (\emph{I.~Ya.~Pomeranchuk}, 1948). The probability of decay can be related to the cross section of annihilation of a free pair.

\SourcePage{414}{419}

The momenta of the electron and positron in positronium are $\sim me^2/\hbar$, i.e.\ small in comparison with $mc$. Therefore in computing the probability of annihilation one can go to the limit of two particles at rest in the origin of coordinates. Let $\bar{\sigma}_{2\gamma}$ be the cross section of two-photon annihilation of a free pair, averaged over spin directions of both particles. In the non-relativistic limit, according to (88.11)\footnote{Formulas (89.1) and (89.7) are written in ordinary units.},
\begin{equation}
\bar{\sigma}_{2\gamma} = \pi\left(\frac{e^2}{mc^2}\right)^{\!2}\frac{c}{v},
\label{eq:89.1}
\end{equation}
where $v$ is the relative velocity of the particles. We obtain the probability of annihilation $\bar{w}_{2\gamma}$ by multiplying $\bar{\sigma}_{2\gamma}$ by the flux density $v|\psi(0)|^2$. Here $\psi(r)$ is the wave function of the ground state of positronium normalised to~1:
\begin{equation}
\psi = \frac{1}{\sqrt{\pi a^3}}\,e^{-r/a}, \quad a = \frac{2\hbar^2}{me^2}
\label{eq:89.2}
\end{equation}
(the Bohr radius of positronium $a$ is twice the Bohr radius of the hydrogen atom because the reduced mass is half that). This probability, however, corresponds to the averaging over spin states of the initial state. But among the four possible spin states of the two-particle system, only one (with total spin~0) is capable of two-photon annihilation. Therefore the average probability of decay $\bar{w}_{2\gamma}$ is related to the probability of decay of parapositronium $w_0$ by the relation $\bar{w}_{2\gamma} = w_0/4$. Thus,
\begin{equation}
w_0 = 4|\psi(0)|^2 (v\bar{\sigma}_{2\gamma})_{v\to 0}.
\label{eq:89.3}
\end{equation}

Substituting the values of the quantities from (89.1) and (89.2), we obtain for the lifetime of parapositronium
\begin{equation}
\tau_0 = \frac{2\hbar}{mc^2\alpha^5} = 1.23 \cdot 10^{-10}~\text{s}.
\label{eq:89.4}
\end{equation}

We draw attention to the fact that the level width $\Gamma_0 = \hbar/\tau_0$ is small compared with the energy:
\[
|E_\text{gr}| = \frac{me^4}{4\hbar^2} = mc^2\,\frac{\alpha^2}{4}.
\]
It is precisely this circumstance that allows one to treat positronium as a quasi-stationary system.

\SourcePage{415}{420}

In an analogous manner we find that the probability of decay of orthopositronium is related to the spin-averaged cross section of three-photon annihilation of a free pair by the relation
\begin{equation}
w_1 = \frac{4}{3}\bar{w}_{3\gamma} = \frac{4}{3}|\psi(0)|^2(v\bar{\sigma}_{3\gamma})_{v\to 0}
\label{eq:89.5}
\end{equation}
($\frac{3}{4}$ is the relative statistical weight of the state with spin~1). Looking ahead, we note that
\begin{equation}
\bar{\sigma}_{3\gamma} = \frac{4(\pi^2-9)}{3v}\,\alpha\left(\frac{e^2}{mc^2}\right)^{\!2}.
\label{eq:89.6}
\end{equation}
Therefore the lifetime of orthopositronium is
\begin{equation}
\tau_1 = \frac{9\pi}{2(\pi^2-9)}\,\frac{\hbar}{mc^2\alpha^6} = 1.4 \cdot 10^{-7}~\text{s}.
\label{eq:89.7}
\end{equation}
The inequality $\Gamma_1 \ll |E_\text{gr}|$ is in this case satisfied, of course, even better than for parapositronium.

Let us compute the cross section of three-photon annihilation of a free pair (\emph{A.~Ore, J.~L.~Powell}, 1949).

According to (64.18) the cross section of the process under consideration in the centre-of-inertia frame is expressed through the squared amplitude by the formula
\begin{equation}
\begin{split}
d\sigma_{3\gamma} &= \frac{(2\pi)^4|M_{fi}|^2}{4I}\,\delta(\mathbf{k}_1 + \mathbf{k}_2 + \mathbf{k}_3)\,\delta(\omega_1 + \omega_2 + \omega_3 - 2m) \times {} \\
&\quad \times \frac{d^3k_1\,d^3k_2\,d^3k_3}{(2\pi)^9 2\omega_1 \cdot 2\omega_2 \cdot 2\omega_3},
\end{split}
\label{eq:89.8}
\end{equation}
where, according to (64.16), $I = 2m\frac{m}{2}v = m^2 v$, where $v$ is the relative velocity of the positron and electron (which we assume to be small); $\mathbf{k}_1$, $\mathbf{k}_2$, $\mathbf{k}_3$ and $\omega_1$, $\omega_2$, $\omega_3$ are the wave vectors and frequencies of the emitted photons; the $\delta$-functions express the laws of conservation of energy and momentum. By virtue of these laws the three frequencies $\omega_1$, $\omega_2$, $\omega_3$ must be representable as the side lengths of a triangle with perimeter~$2m$.

Three-photon annihilation corresponds to the diagram
% [Feynman diagram: Three-photon annihilation. Electron line from p₋ through three vertices emitting photons k₁, k₂, k₃, with positron line from p₊. Plus five more diagrams obtained by permutations of k₁, k₂, k₃]

and five more diagrams obtained by permutation of the photons $k_1$, $k_2$, $k_3$. The corresponding amplitude is written in the form
\begin{equation}
M_{fi} = (4\pi)^{3/2} e_\lambda^{(3)*} e_\mu^{(2)*} e_\nu^{(1)*} \bar{u}(-p_+)Q^{\lambda\mu\nu}u(p_-),
\label{eq:89.9}
\end{equation}

\SourcePage{416}{421}

where
\begin{equation}
Q^{\lambda\mu\nu} = \sum_\text{perm} \gamma^\lambda G(k_3 - p_+)\gamma^\mu G(p_- - k_1)\gamma^\nu,
\label{eq:89.10}
\end{equation}
where the sum is taken over all permutations of the photon numbers 1, 2, 3 together with the simultaneously corresponding permutations of the tensor indices $\lambda\mu\nu$. The squared amplitude, averaged over the polarisations of the electron and positron and summed over the polarisations of the photons:
\begin{equation}
\frac{1}{4}\sum_\text{polar} |M_{fi}|^2 = (4\pi)^3\operatorname{Sp}\!\left\{\rho_+Q^{\lambda\mu\nu}\rho_-\widetilde{Q}_{\lambda\mu\nu}\right\},
\label{eq:89.11}
\end{equation}
where
\[
\rho_- = \frac{1}{2}(\gamma p_- + m), \quad \rho_+ = \frac{1}{2}(\gamma p_+ - m).
\]

The matrices $\widetilde{Q}^{\lambda\mu\nu}$ differ from the matrices $Q^{\lambda\mu\nu}$ by the reversal of the order of factors in each term of the sum. In the case of interest to us, where the velocities of the electron and positron are small, one can set their 3-momenta $\mathbf{p}_-$ and $\mathbf{p}_+$ equal to zero, i.e.\ set $p_- = p_+ = (m, 0)$. Then the electron Green's functions become
\[
G(p_- - k_1) = \frac{\gamma p_- - \gamma k_1 + m}{(p_- - k_1)^2 - m^2} \approx \frac{-\gamma k_1 + m(\gamma^0 + 1)}{-2m\omega_1}
\]
and so on, and the density matrices reduce to
\[
\rho_\pm = \frac{m}{2}(\gamma^0 \pm 1).
\]

Upon multiplying out the terms in (89.11) there arises a large number of terms. However, the number of terms requiring computation can be considerably reduced by making full use of the symmetry with respect to the permutations of the photons. It suffices to multiply out the six terms in $Q^{\lambda\mu\nu}$ (89.10) with any one term in $\widetilde{Q}^{\lambda\mu\nu}$. In the remaining five traces one can also identify certain parts that transform into one another under different permutations of the photons. The products of traces of 4-vectors $p$, $k_1$, $k_2$, $k_3$ arising upon expanding the traces are expressed through the frequencies $\omega_1$, $\omega_2$, $\omega_3$. Since $p = (m, 0)$, we have $pk_1 = m\omega_1$, \dots. The products $k_1 k_2$, \dots\ are determined from the equation of 4-momentum conservation: $2p = k_1 + k_2 + k_3$; rewriting this equation in the form $2p - k_3 = k_1 + k_2$ and squaring, we obtain
\begin{equation}
k_1 k_2 = 2m(m - \omega_3), \ldots
\label{eq:89.12}
\end{equation}

\SourcePage{417}{422}

As a result, after a rather lengthy calculation one obtains
\[
\frac{1}{4}\sum_\text{polar} |M_{fi}|^2 = (4\pi)^3 e^6 \cdot 16\left[\left(\frac{m-\omega_1}{\omega_2\omega_3}\right)^{\!2} + \left(\frac{m-\omega_2}{\omega_1\omega_3}\right)^{\!2} + \left(\frac{m-\omega_3}{\omega_1\omega_2}\right)^{\!2}\right].
\]

Substituting this expression into (89.8), we find the differential cross section of three-photon annihilation:
\begin{equation}
\begin{split}
d\bar{\sigma}_{3\gamma} &= \frac{e^6}{\pi^2 m^2 v}\left[\left(\frac{m-\omega_1}{\omega_2\omega_3}\right)^{\!2} + \left(\frac{m-\omega_2}{\omega_1\omega_3}\right)^{\!2} + \left(\frac{m-\omega_3}{\omega_1\omega_2}\right)^{\!2}\right] \times {} \\
&\quad \times \delta(\mathbf{k}_1 + \mathbf{k}_2 + \mathbf{k}_3)\,\delta(\omega_1 + \omega_2 + \omega_3 - 2m)\,\frac{d^3k_1\,d^3k_2\,d^3k_3}{\omega_1\omega_2\omega_3}.
\end{split}
\label{eq:89.13}
\end{equation}
Here we must still eliminate the $\delta$-functions. The first of them is removed by integration over $d^3k_3$, after which we replace the remaining differentials:
\[
d^3k_1\,d^3k_2 \to 4\pi\omega_1^2\,d\omega_1 \cdot 2\pi\omega_2^2\,d(\cos\theta_{12})\,d\omega_2,
\]
where $\theta_{12}$ is the angle between $\mathbf{k}_1$ and $\mathbf{k}_2$; it is understood that the integration over the directions of $\mathbf{k}_1$ and the azimuth of $\mathbf{k}_2$ relative to $\mathbf{k}_1$ has already been performed. Differentiating the equation
\[
\omega_3 = \sqrt{\omega_1^2 + \omega_2^2 + 2\omega_1\omega_2\cos\theta_{12}},
\]
we find
\[
d\cos\theta_{12} = \frac{\omega_3}{\omega_1\omega_2}\,d\omega_3.
\]
By integrating over $\omega_3$ we eliminate the second $\delta$-function. As a result we obtain the cross section of annihilation with the production of photons having given energies in the form
\begin{equation}
\begin{split}
d\bar{\sigma}_{3\gamma} &= \frac{1}{6}\,\frac{8e^6}{vm^2}\biggl\{\!\left(\frac{m-\omega_1}{\omega_2\omega_3}\right)^{\!2} + \left(\frac{m-\omega_2}{\omega_1\omega_3}\right)^{\!2} + \left(\frac{m-\omega_3}{\omega_1\omega_2}\right)^{\!2}\!\biggr\}\,d\omega_1\,d\omega_2
\end{split}
\label{eq:89.14}
\end{equation}
(with a view to subsequent integration over the frequencies, we have introduced the factor $\frac{1}{6}$ taking into account the identity of the photons (cf.\ the footnote on p.~286)).

Each of the frequencies $\omega_1$, $\omega_2$, $\omega_3$ can take values between 0 and $m$ (the value $m$ is attained by two of the frequencies when the third equals zero). At given $\omega_1$ the frequency $\omega_2$ varies between $m - \omega_1$ and $m$. Integrating (89.14) over $\omega_2$ in these limits, we obtain the spectral distribution of the photons of the decay:
\[
d\bar{\sigma}_{3\gamma} = \frac{8e^6}{3vm^3}\,F(\omega_1)\,d\omega_1,
\]
\[
F(\omega_1) = \frac{\omega_1(m-\omega_1)}{(2m-\omega_1)^2} + \frac{2m-\omega_1}{\omega_1} + \left[\frac{2m(m-\omega_1)}{\omega_1^2} - \frac{2m(m-\omega_1)^2}{(2m-\omega_1)^3}\right]\ln\frac{m-\omega_1}{m}.
\]
The function $F$ monotonically increases from zero at $\omega_1 = 0$ to 1 at $\omega_1 = m$; in Fig.~14 its graph is shown.

% [Figure 14: Graph of F(ω₁) vs ω₁/m, monotonically increasing from 0 to 1]

\SourcePage{418}{423}

The total annihilation cross section is obtained by integrating (89.14) over both frequencies:
\[
\bar{\sigma}_{3\gamma} = \frac{4e^6}{3vm^2} \cdot 3 \int_0^m\!\int_{m-\omega_1}^m \frac{(\omega_1+\omega_2-m)^2}{\omega_1^2\omega_2^2}\,d\omega_1\,d\omega_2.
\]
The double integral here equals $(\pi^2 - 9)/3$, and we arrive at the formula (89.6) given above.
